Во второй главе рассматривается архитектура разработанного программного комплекса. Основной акцент сделан на вычислительном модуле \texttt{compute}, поскольку именно он является центральным результатом дипломной работы. Серверный слой на gRPC и клиент визуализации на Flutter рассматриваются как вспомогательные компоненты, разработанные для удобства экспериментов, удалённого запуска и анализа результатов.

\subsection{Роль вычислительного модуля в структуре проекта}

Целью дипломной работы является реализация параллельных алгоритмов интерполяции на основе разреженных сеток с квадратичным базисом. Этой цели непосредственно соответствует модуль \texttt{compute}, в котором сосредоточены все ключевые математические и алгоритмические решения: представление узлов, вычисление базисных функций, построение адаптивной сетки, рекурсивная интерполяция, композиция операторов и поддержка режимов последовательного и параллельного построения.

Сервисный слой и пользовательский интерфейс не определяют математическое содержание работы, а лишь расширяют область практического применения созданного ядра. По этой причине архитектура программного комплекса должна рассматриваться прежде всего как архитектура вычислительного модуля, дополненная средствами визуализации и сетевого доступа.

\subsection{Требования к вычислительному модулю}

К модулю \texttt{compute} были предъявлены следующие требования:
\begin{enumerate}
    \item поддержка многомерных входных данных и векторных выходных значений;
    \item поддержка линейного и квадратичного базисов;
    \item адаптивное уточнение сетки на основе коэффициентов избытка;
    \item возможность учёта якорных точек для ручного управления уточнением;
    \item поддержка последовательного и параллельного режимов построения;
    \item возможность жёсткого ограничения числа узлов;
    \item поддержка композиции отображений и последовательных итераций оператора;
    \item совместимость с задачами прямой интерполяции и численного решения дифференциальных уравнений.
\end{enumerate}

Именно выполнение этих требований определяет архитектурные решения, рассмотренные далее.

\subsection{Центральный тип AdaptiveSparseGrid}

Ключевой структурой вычислительного модуля является тип \texttt{AdaptiveSparseGrid}. Он объединяет в одном объекте:
\begin{enumerate}
    \item границы области определения;
    \item тип используемого базиса;
    \item размерность входного и выходного пространств;
    \item набор входных точек построения;
    \item словарь всех узлов сетки.
\end{enumerate}

Такое проектирование позволяет трактовать объект сетки одновременно как:
\begin{enumerate}
    \item контейнер данных, описывающий уже построенный интерполянт;
    \item исполняемый оператор, умеющий вычислять приближённое значение в новой точке;
    \item основу для построения следующей итерации или композиции отображений.
\end{enumerate}

Это архитектурное свойство особенно важно для задач моделирования динамических систем, где один и тот же объект должен не только хранить аппроксимацию, но и участвовать в последующих шагах вычисления.

\subsection{Проектирование базовых структур данных}

Для реализации метода были введены следующие типы.

\textbf{Тип \texttt{Point}.} Представляет собой вектор вещественных координат. Он используется как для входных точек, так и для значений функции и коэффициентов избытка. Использование единого векторного представления позволяет поддерживать как скалярные, так и векторные отображения без дублирования алгоритмов.

\textbf{Тип \texttt{gridKey}.} Определяет положение узла в иерархической структуре через массивы уровней и индексов. Такой ключ задаёт не только идентификатор узла, но и его геометрический смысл: уровень отвечает за масштаб, индекс --- за положение центра.

\textbf{Тип \texttt{node}.} Содержит ключ, координаты центра в единичной области, коэффициент избытка и признак наличия потомков. Через эту структуру материализуется математическая модель адаптивной сетки.

\textbf{Тип \texttt{entryPoint}.} Описывает стартовую конфигурацию построения. Входные точки играют важную роль, поскольку именно от них запускаются независимые локальные ветви построения, что напрямую связано с возможностью распараллеливания.

В совокупности эти структуры образуют минимальное, но достаточное ядро вычислительной модели.

\subsection{Архитектура построения сетки}

С точки зрения архитектуры модуль \texttt{compute} организован вокруг процедуры построения, которая разделяется на два уровня.

На первом уровне формируются стартовые конфигурации уровней и индексов, соответствующие различным входным точкам. На втором уровне для каждой такой конфигурации запускается локальная задача адаптивного построения. Локальная задача:
\begin{enumerate}
    \item создаёт начальный узел;
    \item вычисляет коэффициент избытка;
    \item проверяет критерии продолжения;
    \item при необходимости порождает дочерние узлы;
    \item рекурсивно или итеративно расширяет сетку до выполнения ограничений.
\end{enumerate}

Это решение позволяет отделить глобальную организацию стартовых конфигураций от локальной логики уточнения. Такое разделение удобно и математически, и программно: оно делает код более модульным и упрощает последующую параллелизацию.

\subsection{Архитектура поддержки parallel и sequential build}

В модуле предусмотрены два режима работы: \texttt{BuildTypeSequential} и \texttt{BuildTypeParallel}. Их поддержка реализована не как два разных алгоритма, а как два режима исполнения одной и той же логики построения.

В последовательном режиме все локальные задачи выполняются по очереди в одном потоке. Этот режим удобен для проверки корректности, профилирования и сопоставления с параллельным вариантом.

В параллельном режиме независимые стартовые конфигурации обрабатываются одновременно. Архитектурное достоинство такого подхода состоит в том, что конкурентность не «вшивается» в математические операции над узлом, а накладывается поверх уже выделенных независимых задач. Благодаря этому ядро сохраняет логическую простоту, а параллелизация не разрушает его структуру.

\subsection{Поддержка якорных точек как архитектурное расширение}

В ходе разработки была выявлена проблема, когда адаптивный критерий, опирающийся только на коэффициент избытка в центре узла, может преждевременно остановить уточнение. Для решения этой проблемы в архитектуру вычислительного модуля был включён дополнительный механизм якорных точек.

С архитектурной точки зрения якорные точки представляют собой внешний набор контрольных точек, передаваемых в конструктор сетки. Они не изменяют базовую формулу интерполяции, но влияют на принятие решения о продолжении построения. Это хороший пример расширения алгоритма без изменения его фундаментальной структуры: в ядро добавляется дополнительный слой условий, не нарушающий основной логики вычисления коэффициентов и рекурсивного обхода.

\subsection{Жёсткое ограничение числа узлов как эксплуатационный параметр}

Важным архитектурным решением является поддержка параметра \texttt{maxNodesInGrid}. Он задаёт верхнюю границу числа узлов, после которой построение должно завершиться принудительно. Наличие этого ограничения позволяет использовать вычислительный модуль не только в исследовательском режиме, но и в условиях ограниченных ресурсов.

С практической точки зрения это означает, что пользователь может выбрать один из двух режимов управления сложностью:
\begin{enumerate}
    \item через точность и глубину уточнения;
    \item через жёсткий лимит на размер сетки.
\end{enumerate}

Такое решение особенно важно для серверного использования, где необходимо контролировать время ответа и объём данных, передаваемых клиенту.

\subsection{Композиция отображений и метод следующей итерации}

Одной из сильнейших сторон архитектуры модуля \texttt{compute} является поддержка композиции отображений через метод \texttt{MakeNextIteration}. В отличие от обычного вычислительного объекта, сетка в этом случае используется как промежуточный оператор: её результат подаётся на вход следующей функции, после чего строится новая аппроксимация.

С архитектурной точки зрения это превращает \texttt{AdaptiveSparseGrid} в универсальный вычислительный кирпич, пригодный для последовательного моделирования шагов эволюции. Именно поэтому модуль \texttt{compute} подходит не только для статической интерполяции, но и для задач динамики.

\subsection{Поддержка дифференциальных уравнений на уровне архитектуры}

Для работы с дифференциальными уравнениями используются два концептуально разных режима.
\begin{enumerate}
    \item \textbf{Параметр $t$ как интервальная неопределённость.} В этом случае время включается как дополнительная координата входного пространства. Тогда аппроксимируется отображение, зависящее как от начального состояния, так и от значения времени.
    \item \textbf{Итеративный режим через \texttt{MakeNextIteration}.} В этом случае строится одношаговый оператор перехода, после чего он последовательно композиционно применяется несколько раз.
\end{enumerate}

Архитектурная ценность этой двойной схемы состоит в том, что один и тот же вычислительный модуль способен поддерживать два разных численных сценария без изменения своего интерфейса. Это также делает возможным содержательное сравнение их производительности на уровне бенчмарков.

\subsection{Вспомогательный серверный слой}

Слой \texttt{internal/server} не является центральным результатом дипломной работы, но играет важную вспомогательную роль. Его задача --- предоставить удобный удалённый интерфейс к вычислительному модулю. Сервер:
\begin{enumerate}
    \item принимает параметры задачи;
    \item интерпретирует пользовательские формулы;
    \item при необходимости решает дифференциальное уравнение численным методом;
    \item передаёт полученное отображение в \texttt{compute};
    \item сериализует построенную сетку в protobuf-ответ.
\end{enumerate}

Таким образом, серверный слой следует рассматривать как адаптер между пользовательским запросом и вычислительным ядром, а не как самостоятельный предмет исследования.

\subsection{Клиентская визуализация как вспомогательный инструмент}

Аналогичную роль играет и Flutter-приложение. Оно предназначено для:
\begin{enumerate}
    \item ввода параметров задачи;
    \item запуска расчётов без работы из командной строки;
    \item визуального анализа распределения узлов;
    \item демонстрации результатов в наглядной форме.
\end{enumerate}

Клиентская часть важна для практического применения проекта, однако архитектурный центр диплома находится именно в модуле \texttt{compute}. Flutter-интерфейс выступает как средство визуализации результатов работы этого модуля.

\subsection{Наблюдаемость и эксплуатационные характеристики}

Несмотря на вспомогательную роль серверного слоя, в проекте предусмотрены средства эксплуатационного мониторинга: метрики Prometheus, контроль длительности запроса, числа узлов и размера ответа. Эти механизмы полезны прежде всего для оценки того, как параметры вычислительного модуля отражаются на внешних характеристиках системы.

Иначе говоря, наблюдаемость сервиса введена не ради самой серверной архитектуры, а как средство исследования поведения алгоритмов, реализованных в \texttt{compute}.

\subsection*{Выводы по главе}
\addcontentsline{toc}{subsection}{Выводы по главе}

Во второй главе была рассмотрена архитектура программного комплекса с акцентом на вычислительный модуль \texttt{compute}, являющийся главным результатом дипломной работы. Показано, что именно в этом модуле сосредоточены все ключевые математические и алгоритмические решения: представление узлов, адаптивное уточнение, поддержка якорных точек, ограничение числа узлов, последовательный и параллельный режимы построения, а также композиция операторов и работа с дифференциальными уравнениями. Серверный слой на gRPC и клиент визуализации на Flutter рассмотрены как вспомогательные средства удалённого доступа и анализа результатов. Такая архитектура подчёркивает, что ядром проекта является именно реализация параллельных алгоритмов интерполяции на основе разреженных сеток.
